% =========================================================
% 101.introduction_revisada.tex

\section{Introdução}
\label{sec:introducao}

Nematoides fitoparasitas reduzem a produtividade agrícola ao colonizar tecidos vegetais e comprometer a absorção de água e nutrientes. A identificação da espécie e a estimativa da população orientam o manejo, pois cada grupo possui ciclo biológico, faixa de hospedeiros e limiar de dano próprios \cite{agarwal2023}.

O nematoide de cisto da soja, \textit{Heterodera glycines} Ichinohe, infecta as raízes e estabelece sítios de alimentação no cilindro vascular. A infecção prejudica o sistema radicular, reduz o porte da planta e causa perdas de rendimento. A ausência de sintomas aéreos específicos dificulta a detecção por inspeção visual \cite{he2023}.

A persistência dos cistos no solo e a adaptação das populações dificultam o manejo. Cultivares resistentes e rotação de culturas ajudam a controlar o patógeno, mas o uso contínuo da mesma fonte de resistência seleciona populações virulentas. A durabilidade da resistência depende da genética do hospedeiro, da virulência do nematoide e do número de gerações por ciclo de cultivo \cite{mccarville2024}. A escolha do manejo requer dados sobre a presença do patógeno e as características da área infestada.

O diagnóstico utiliza amostras de solo ou raízes, extração dos organismos e identificação morfológica ou molecular. A identificação morfológica exige microscopia e conhecimento taxonômico. Os métodos moleculares oferecem especificidade, mas requerem preparo da amostra, reagentes e equipamentos. O volume de amostras nos laboratórios aumenta o prazo para emissão dos resultados \cite{agarwal2023}.

A qualidade do diagnóstico depende da representatividade, identificação, conservação e transporte das amostras. Falhas nessas etapas comprometem o material. A ocorrência e a densidade dos nematoides variam entre pontos da mesma área, o que exige planos de amostragem compatíveis com a heterogeneidade dos talhões \cite{kessler2023}. O aumento do número de pontos melhora a resolução espacial, mas eleva os custos com recipientes, registros, transporte e análises. Uma ferramenta portátil de triagem pode indicar os pontos que exigem confirmação nematológica.

Esse sistema não substitui o diagnóstico laboratorial. Sua função é gerar indícios sobre a amostra, apoiar a seleção de pontos e ampliar a cobertura do monitoramento. Os resultados podem orientar coletas, acompanhar áreas sob manejo e selecionar locais para contagem ou identificação molecular. A resposta sensorial não representa identificação taxonômica do patógeno.

Sensores de umidade, temperatura, pressão e propriedades físico-químicas do solo descrevem o ambiente da cultura, mas apresentam baixa especificidade para \textit{H. glycines}. Déficit hídrico, compactação, atividade microbiana, nutrição e outros agentes bióticos podem gerar respostas semelhantes. Os compostos voláteis do sistema solo--planta fornecem outra fonte de informação, pois alterações fisiológicas e metabólicas modificam a composição do \textit{headspace}.

Narizes eletrônicos empregam matrizes de sensores com seletividade parcial. Cada sensor responde a grupos de gases, e o conjunto das respostas forma uma assinatura multivariada. A classificação depende da matriz sensorial, do protocolo de amostragem, do tratamento dos sinais e do modelo computacional. Sistemas MOS baseados em perfis de compostos orgânicos voláteis reconheceram danos causados por insetos em arroz \cite{li2025}. Matrizes MOS combinadas com aprendizado de máquina detectaram brusone em plantas sem sintomas aparentes \cite{wang2025}. Esses estudos mostram que o conjunto de sinais contém informação sobre o estado fitossanitário, sem exigir a identificação de cada composto.

Na soja, assinaturas gasosas respondem às condições fisiológicas e à origem das amostras. \citeonline{herrmann2024} utilizaram nariz eletrônico e aprendizado de máquina para distinguir plantas sob estresse hídrico e períodos de aquisição. \citeonline{sun2024} combinaram nariz eletrônico e rede neural com atenção de canais para reconhecer amostras de diferentes origens. Esses resultados sustentam o uso de padrões gasosos, mas indicam fatores de confusão, como umidade, horário, ambiente, origem e condição fitossanitária.

Umidade e temperatura afetam a superfície dos sensores semicondutores de óxido metálico. Pressão e fluxo alteram o transporte dos gases até a câmara. O teor de água no solo interfere na liberação e na retenção dos compostos do sistema solo--planta. Em um estudo sobre a detecção de vírus em pimenta, o controle do ruído associado às condições de aquisição aumentou o desempenho do nariz eletrônico e evidenciou a importância do pré-processamento \cite{agustika2025}. A avaliação em faixas distintas de umidade permite medir a estabilidade da separação entre as classes.

O modo de coleta afeta os sinais. A coleta a vácuo conduz o \textit{headspace} até os sensores e altera a amplitude, o tempo de exposição e o perfil temporal da resposta. A coleta passiva depende da difusão dos gases. A comparação entre esses modos permite medir a influência do protocolo sobre o classificador. Sem esse controle, o modelo pode associar características da coleta ao rótulo biológico.

As respostas do nariz eletrônico formam séries temporais multivariadas com transientes de subida, regiões de resposta e variações da interação entre gases e sensores. Médias e valores máximos reduzem a sequência a poucos descritores e descartam relações temporais. Redes convolucionais temporais usam convoluções dilatadas e conexões residuais para ampliar o campo receptivo e extrair padrões ao longo do tempo. Estudos com narizes eletrônicos aplicaram TCNs à aprendizagem de relações temporais, à combinação de informações de tempo e frequência e à classificação de gases em ambientes abertos \cite{ma2024tftcn,ma2024msetcn}.

A formação dos conjuntos de treinamento e teste deve considerar a estrutura dos dados. Leituras da mesma coleta compartilham vaso, planta, dia, ambiente e protocolo. A divisão aleatória dessas leituras pode distribuir fragmentos da mesma unidade experimental entre os dois conjuntos e causar vazamento de dados. Nesse caso, o modelo reconhece a coleta em vez da condição biológica. A separação por coleta preserva a independência dos grupos e representa o uso do modelo em novas aquisições.

Este trabalho investiga a discriminação entre plantas de soja do grupo inoculado e do grupo controle por meio de um nariz eletrônico de baixo custo. O sistema reúne seis sensores semicondutores (MQ-2, MQ-3, MQ-7, MQ-8, MQ-135 e MQ-138) e mede umidade do solo, temperatura e pressão. Algoritmos de aprendizado de máquina processam os sinais organizados em janelas temporais. A divisão entre treinamento e teste ocorre por coleta.

A análise abrange três cenários: solo seco com coleta a vácuo, solo úmido com coleta a vácuo e coleta passiva sem estratificação por umidade. A avaliação utiliza acurácia, acurácia balanceada e F1-macro. A contribuição tecnológica reúne aquisição sob condições físicas distintas, compensação de variáveis ambientais, modelagem temporal e validação agrupada. O estudo identifica as condições em que a assinatura sensorial permite separar as classes e os casos em que a umidade e o modo de coleta prejudicam essa separação. Os resultados constituem uma prova de conceito para triagem em ambiente controlado. A confirmação da infestação cabe aos métodos nematológicos ou moleculares.
